\documentclass[11pt]{article}
\usepackage{mathblatt}


\begin{document}
\pfheftstil
\pfheftkopf{Wachstum}{P10}
\pfrueckblick
\begin{pfaufg}{}{1.}{}
Ergänze die Tabelle Änderung | Faktor.
\par\smallskip\noindent \begingroup\renewcommand{\arraystretch}{1.3}\begin{tabular}{|c|c|}\hline \textbf{Änderung} & \textbf{Faktor} \\ \hline 1,04 & \rule{0pt}{3ex}\hspace{1.6cm} \\ \hline 0,87 & \rule{0pt}{3ex}\hspace{1.6cm} \\ \hline 1,03 & \rule{0pt}{3ex}\hspace{1.6cm} \\ \hline 1,019 & \rule{0pt}{3ex}\hspace{1.6cm} \\ \hline\end{tabular}\endgroup\par
\fusshilfe{\mbox{1:~1,04; 0,87; 1,03; 1,019}}
\end{pfaufg}
\begin{pfaufg}{}{2.}{}
80 min sollen auf 16 Kästchen passen: Wie viele Minuten je Kästchen? 1000 hPa auf 20 Kästchen: wie viele hPa je Kästchen?
\fusshilfe{\mbox{2:~5 min; 50 hPa}}
\end{pfaufg}
\begin{pfaufg}{}{3.}{}
Berechne: 112 : 80; 648 000 : 600 000; 1,04².
\rechenraster{1}
\fusshilfe{\mbox{3:~1,4; 1,08; 1,0816}}
\end{pfaufg}
\pfstufekopf[stabelleergnzen]{Tabelle ergänzen}{in 1 der letzten 5 Prüfungen}
\pfgruppe{}
\begin{pfaufg}{}{4.}{}
Die Tabelle zeigt Werte. Kreuze an, was von Wert zu Wert gleich bleibt. \wertetabelle[100,110,121]{t}{\text{Wert}}{0,1,2}
\pfkreuzzeile{die Differenz}\pfkreuzzeile{der Quotient}
\fusshilfe{\mbox{4:~der Quotient}}
\end{pfaufg}
\pfgruppe{}
\begin{pfaufg}{NRW ’23}{5.}{}
In einem Bienenvolk leben 308 Varroa-Milben. Ihre Zahl verdoppelt sich etwa alle vier Wochen. Ergänze die Tabelle.
\par\smallskip\noindent \sachtabelle{l}{Zeit in Wochen 0, 4, 8}{Anzahl der Milben 308, leer, leer}\par
\fusshilfe{\mbox{5:~616 und 1 232}}
\end{pfaufg}
\pfgruppe{}
\begin{pfaufg}{P10 ’20}{6.}{}
Laut einer Studie aus dem Jahr 2019 verbrauchte jeder Mensch im Jahr 2019 im Schnitt 54 Tausend Liter Wasser; der Verbrauch je Person soll in Zukunft jedes Jahr um 3 \% wachsen. Ergänze die Tabelle mit diesen Angaben.
\par\smallskip\noindent \par\medskip\begingroup\centering\renewcommand{\arraystretch}{1.25}\begin{tabular}{|p{4.5cm}|p{1.6cm}|p{1.6cm}|p{1.6cm}|p{1.6cm}|}\hline Jahr & 2019 & 2020 & 2021 & 2022 \\ \hline Wasserverbrauch je Person in Tausend Litern & \leerzelle & 55,6 & 57,3 & \leerzelle \\ \hline\end{tabular}\par\endgroup\medskip\par
\fusshilfe{\mbox{6:~2019: 54; 2022: 57,3 · 1,03 = 59,019 $\approx$ 59,0}}
\end{pfaufg}
\pfgruppe{}
\begin{pfaufg}{P10 ’16}{7.}{}
Um Krankheiten der Schilddrüse zu erkennen, spritzt man dem Patienten eine Flüssigkeit, die eine kleine Menge radioaktives Technetium enthält. Die Masse des Technetiums im Körper wird jede Stunde um etwa 11 \% kleiner. Ergänze die beiden Lücken in der Tabelle: die Masse nach 1 Stunde und die Zeit, zu der noch 3,14 mg vorhanden sind.
\par\smallskip\noindent \sachtabelle{lccccccc}{Zeit in h & 0 & 1 & 2 & \leerzelle & 5 & 6 & 8}{Masse in mg & 5,00 & \leerzelle & 3,96 & 3,14 & 2,79 & 2,48 & 1,97}\par
\pfzwfrage{Gib den Faktor an, mit dem die Masse jede Stunde multipliziert wird.}
\pfzwfrage{Multipliziere 5,00 mg mit diesem Faktor.}
\pfzwfrage{Rechne von 2 h aus Stunde für Stunde weiter, bis 3,14 mg erreicht sind.}
\fusshilfe{\mbox{7:~4,45 mg; 4 h}}
\end{pfaufg}
\pfgruppe{}
\begin{pfaufg}{P10 ’17}{8.}{}
Auf Meereshöhe (0 km) misst man einen Luftdruck von ungefähr 1000 hPa (Hektopascal). Je höher man steigt, desto kleiner wird der Luftdruck: mit jedem Kilometer Höhe um 13 \%. Trag die beiden fehlenden Werte in die Tabelle ein.
\par\smallskip\noindent \sachtabelle{lccccccc}{Höhe in km & 0 & 1 & 2 & 3 & 5 & 8 & 10}{Luftdruck in hPa & 1000 & 870 & \leerzelle & 659 & 498 & \leerzelle & 248}\par
\pfzwfrage{Gib den Faktor an, mit dem der Luftdruck je Kilometer multipliziert wird.}
\fusshilfe{\mbox{8:~756,9 $\approx$ 757 hPa; $\approx$ 328 hPa}}
\end{pfaufg}
\pfgruppe{}
\begin{pfaufg}{P10 ’19}{9.}{}
Ein Ferkel wiegt 10 kg. Aus Erfahrung weiß man, dass seine Masse jede Woche um etwa 4 \% zunimmt. Zu Beginn jeder Woche wird es gewogen. Trag die drei fehlenden Werte in die Tabelle ein.
\par\smallskip\noindent \sachtabelle{lccccc}{Zeit in Wochen & 0 & 1 & 2 & … & 5}{Masse in kg & \leerzelle & 10,400 & \leerzelle & – & \leerzelle}\par
\fusshilfe{\mbox{9:~10 kg; 10,816 kg; $\approx$ 12,167 kg}}
\end{pfaufg}

\pfstufekopf[spunktedarstellen]{Punkte darstellen}{in 1 der letzten 5 Prüfungen}
\pfgruppe{}
\begin{pfaufg}{}{10.}{}
Der Zahlenstrahl ist in Hunderterschritten beschriftet. Welche Zahl gehört zum Punkt $A$?
\par\smallskip\noindent \zahlenstrahl[xmin=0,xmax=600,xstep=100,karo=1]{350/A}\par
\par\noindent \leerfeld \par
\fusshilfe{\mbox{10:~$350$}}
\end{pfaufg}
\pfgruppe{}
\begin{pfaufg}{P10 ’25}{11.}{}
Man beobachtet einen Bakterienstamm, der exponentiell wächst. Die Tabelle zeigt die Zahl der Bakterien in den ersten 4 Stunden. Stelle das Wachstum der Bakterien im Koordinatensystem dar.
\par\smallskip\noindent \begin{tikzpicture}\begin{axis}[xmin=0,xmax=4,ymin=0,ymax=400,tick label style={font=\scriptsize},clip=true,/pgf/number format/use comma,axis line style={-{Stealth[length=2mm]}},every axis plot/.append style={line width=0.9pt},width=11.5cm,height=7.5cm,grid=both,grid style={mbgitter,line width=0.3pt},minor tick num=1,axis lines=left,xlabel={Zeit in Stunden},ylabel={Anzahl der Bakterien},label style={font=\small},scaled ticks=false,/pgf/number format/1000 sep={\,}]\end{axis}\end{tikzpicture}\par
\rechenraster{1}
\end{pfaufg}
\begin{pfaufg}{P10 ’21}{12.}{}
Eine 40 cm hohe Kerze brennt ab. Gegeben: Tabelle Zeit in Minuten 0, 10, 20, 30, 40, 80; Höhe in cm 40, 38, 36, 34, 32, 24. Stelle diese Werte als Graph im Koordinatensystem dar. Trag dazu zuerst die fehlende Einteilung der Achsen ein.
\par\smallskip\noindent \begin{tikzpicture}[x=0.450cm,y=0.450cm]\draw[mbgitter,line width=0.3pt,step=1] (0,0) grid (26,24);\draw[line width=0.7pt,-{Stealth[length=2mm]}] (1,1) -- (25.7,1);\draw[line width=0.7pt,-{Stealth[length=2mm]}] (1,1) -- (1,23.7);\node[font=\small,anchor=north east] at (25.7,0.8) {Zeit in Minuten};\node[font=\small,anchor=south west] at (1.3,23.7) {Höhe in cm};\end{tikzpicture}\par
\pfzwfrage{Lege fest, wie viele Minuten ein Kästchen auf der Rechtsachse bedeutet.}
\pfzwfrage{Lege fest, wie viele Zentimeter ein Kästchen auf der Hochachse bedeutet.}
\rechenraster{2}
\fusshilfe{\mbox{12:~2 cm}}
\end{pfaufg}
\pfgruppe{}
\begin{pfaufg}{P10 ’16}{\llap{\textasteriskcentered\,}13.}{}
Um Krankheiten der Schilddrüse zu erkennen, spritzt man dem Patienten eine Flüssigkeit, die eine kleine Menge radioaktives Technetium enthält. Die Masse des Technetiums im Körper wird jede Stunde um etwa 11 \% kleiner. Beschrifte die Achsen des Koordinatensystems passend zur Tabelle und zeichne ein, wie die Masse abnimmt. Gegeben: Tabelle 0 h: 5,00 mg; 1 h: 4,45 mg; 2 h: 3,96 mg; 4 h: 3,14 mg; 5 h: 2,79 mg; 6 h: 2,48 mg; 8 h: 1,97 mg.
\par\smallskip\noindent \begin{tikzpicture}[x=0.450cm,y=0.450cm]\draw[mbgitter,line width=0.3pt,step=1] (0,0) grid (19,18);\draw[line width=0.7pt,-{Stealth[length=2mm]}] (1,2) -- (18.7,2);\draw[line width=0.7pt,-{Stealth[length=2mm]}] (1,2) -- (1,17.7);\draw[line width=0.5pt] (3,1.75) -- (3,2.25);\draw[line width=0.5pt] (5,1.75) -- (5,2.25);\draw[line width=0.5pt] (7,1.75) -- (7,2.25);\draw[line width=0.5pt] (9,1.75) -- (9,2.25);\draw[line width=0.5pt] (11,1.75) -- (11,2.25);\draw[line width=0.5pt] (13,1.75) -- (13,2.25);\draw[line width=0.5pt] (15,1.75) -- (15,2.25);\draw[line width=0.5pt] (17,1.75) -- (17,2.25);\draw[line width=0.5pt] (0.75,4) -- (1.25,4);\draw[line width=0.5pt] (0.75,6) -- (1.25,6);\draw[line width=0.5pt] (0.75,8) -- (1.25,8);\draw[line width=0.5pt] (0.75,10) -- (1.25,10);\draw[line width=0.5pt] (0.75,12) -- (1.25,12);\draw[line width=0.5pt] (0.75,14) -- (1.25,14);\draw[line width=0.5pt] (0.75,16) -- (1.25,16);\end{tikzpicture}\par
\pfzwfrage{Wähle, welche Größe auf die Rechtsachse und welche auf die Hochachse kommt.}
\rechenraster{2}
\fusshilfe{\mbox{13:~y-Achse Masse in mg, 1 mg je Teilstrich}}
\end{pfaufg}
\pfgruppe{}
\begin{pfaufg}{P10 ’17}{14.}{}
Auf Meereshöhe (0 km) misst man einen Luftdruck von ungefähr 1000 hPa (Hektopascal). Je höher man steigt, desto kleiner wird der Luftdruck: mit jedem Kilometer Höhe um 13 \%. Vervollständige das Koordinatensystem und zeichne den Luftdruck in Abhängigkeit von der Höhe. Gegeben: 0 km: 1000 hPa; 1 km: 870 hPa; 2 km: 757 hPa; 3 km: 659 hPa; 5 km: 498 hPa; 8 km: 328 hPa; 10 km: 248 hPa.
\par\smallskip\noindent \begin{tikzpicture}[x=0.450cm,y=0.450cm]\draw[mbgitter,line width=0.3pt,step=1] (0,0) grid (24,24);\draw[line width=0.7pt,-{Stealth[length=2mm]}] (2,2) -- (23.7,2);\draw[line width=0.7pt,-{Stealth[length=2mm]}] (2,2) -- (2,23.7);\end{tikzpicture}\par
\rechenraster{2}
\fusshilfe{\mbox{14:~x-Achse Höhe in km, 2 Kästchen je km}}
\end{pfaufg}

\pfstufekopf[sfaktorbestimmengleic]{Faktor bestimmen, Gleichung aufstellen}{in 2 der letzten 5 Prüfungen}
\pfgruppe{}
\begin{pfaufg}{}{15.}{}
Ein Wert wird jedes Jahr mit $1{,}35$ multipliziert. Kreuze an, ob der Wert mehr oder weniger wird. 
\pfkreuzzeile{mehr}\pfkreuzzeile{weniger}
\fusshilfe{\mbox{15:~der Faktor ist größer als eins}}
\end{pfaufg}
\begin{pfaufg}{}{16.}{}
Berechne $1{,}04^5$. Runde auf vier Stellen nach dem Komma.
\par\noindent $\approx$ \leerfeld \par
\fusshilfe{\mbox{16:~$\approx 1{,}2167$}}
\end{pfaufg}
\pfgruppe{}
\begin{pfaufg}{P10 ’17}{\llap{\textasteriskcentered\,}17.}{}
Auf Meereshöhe (0 km) misst man einen Luftdruck von ungefähr 1000 hPa (Hektopascal). Je höher man steigt, desto kleiner wird der Luftdruck: mit jedem Kilometer Höhe um 13 \%. Die Abnahme des Luftdrucks lässt sich mit einer Gleichung beschreiben ($x$: Höhe in km, $y$: Luftdruck in hPa). Kreuze die richtige an:
\pfkreuzzeile{$y = 0{,}87 \cdot x$}\pfkreuzzeile{$y = 1000 \cdot 0{,}87^x$}\pfkreuzzeile{$y = x^{1{,}13}$}\pfkreuzzeile{$y = 1000 \cdot 1{,}13^x$}
\par\smallskip\noindent \sachtabelle{lccccccc}{Höhe in km & 0 & 1 & 2 & 3 & 5 & 8 & 10}{Luftdruck in hPa & 1000 & 870 & \leerzelle & 659 & 498 & \leerzelle & 248}\par
\fusshilfe{\mbox{17:~y = 1000 · $0{,}87^{x}$}}
\end{pfaufg}
\pfgruppe{}
\begin{pfaufg}{P10 ’18}{18.}{}
Ein junges Berliner Unternehmen zeigt in einer Tabelle, wie sein Umsatz gewachsen ist: 2013: 600 000 €, 2014: 648 000 €, 2015: 699 840 €, 2016: 755 827 €, 2017: 816 293 €. Es behauptet, der Umsatz wachse jedes Jahr um 8 \%. Zeige, dass das für 2014 stimmt.
\par\smallskip\noindent \sachtabelle{lrrrrr}{Jahr & 2013 & 2014 & 2015 & 2016 & 2017}{Umsatz in € & 600 000 & 648 000 & 699 840 & 755 827 & 816 293}\par
\fusshilfe{\mbox{18:~600 000 € · 1,08 = 648 000 €}}
\end{pfaufg}
\begin{pfaufg}{P10 ’19}{\llap{\textasteriskcentered\,}19.}{}
Ein Ferkel wiegt 10 kg. Aus Erfahrung weiß man, dass seine Masse jede Woche um etwa 4 \% zunimmt. Zu Beginn jeder Woche wird es gewogen. Die Bäuerin Wilma berechnet die Masse des Ferkels nach $x$ Wochen mit $f(x) = 10 \cdot 1{,}04^x$. Schreib auf, was $10$, $1{,}04$ und $x$ in dieser Gleichung bedeuten.
\par\smallskip\noindent \sachtabelle{lc}{Teil der Gleichung & Bedeutung}{10 & \leerzelle\\ 1,04 & \leerzelle\\ x & \leerzelle}\par
\fusshilfe{\mbox{19:~1,04: Wachstumsfaktor, +4 \% je Woche}}
\end{pfaufg}
\pfgruppe{}
\begin{pfaufg}{P10 ’17}{20.}{}
Auf Meereshöhe (0 km) misst man einen Luftdruck von ungefähr 1000 hPa (Hektopascal). Je höher man steigt, desto kleiner wird der Luftdruck: mit jedem Kilometer Höhe um 13 \%. Eine Bergsteigerin misst 573 hPa. Entscheide, ob sie ungefähr 4 km hoch sein kann, und begründe.
\par\smallskip\noindent \sachtabelle{lccccccc}{Höhe in km & 0 & 1 & 2 & 3 & 5 & 8 & 10}{Luftdruck in hPa & 1000 & 870 & \leerzelle & 659 & 498 & \leerzelle & 248}\par
\fusshilfe{\mbox{20:~ja, 1000 · 0,87$^{4}$ $\approx$ 573 hPa}}
\end{pfaufg}
\pfgruppe{}
\begin{pfaufg}{P10 ’16}{\llap{\textasteriskcentered\,}21.}{}
Um Krankheiten der Schilddrüse zu erkennen, spritzt man dem Patienten eine Flüssigkeit, die eine kleine Menge radioaktives Technetium enthält. Die Masse des Technetiums im Körper wird jede Stunde um etwa 11 \% kleiner. Berechne, wie viele Milligramm Technetium 24 Stunden nach der Spritze noch im Körper sind.
\par\smallskip\noindent \sachtabelle{lccccccc}{Zeit in h & 0 & 1 & 2 & \leerzelle & 5 & 6 & 8}{Masse in mg & 5,00 & \leerzelle & 3,96 & 3,14 & 2,79 & 2,48 & 1,97}\par
\pfzwfrage{Schreib den Term für die Masse nach $t$ Stunden auf.}
\fusshilfe{\mbox{21:~$\approx$ 0,31 mg}}
\end{pfaufg}
\pfgruppe{}
\begin{pfaufg}{P10 ’20}{\llap{\textasteriskcentered\,}22.}{}
Nach einer Studie lag der Wasserverbrauch 2019 bei durchschnittlich 54 Tausend Litern pro Person und Jahr; er soll jedes Jahr um 3 \% steigen. Diese Vorhersage beschreibt die Gleichung $y = 54\cdot 1{,}03^{x}$ ($x$: Jahre ab 2019, $y$: Verbrauch pro Person in Tausend Litern). Eine zweite Studie erwartet, dass der Verbrauch pro Person von 2019 bis 2033 insgesamt um etwa 40 \% zunimmt. Untersuche, ob beide Studien für 2033 denselben Verbrauch vorhersagen.
\pfzwfrage{Berechne mit der Gleichung den Verbrauch pro Person im Jahr 2033.}
\rechenraster{3}
\fusshilfe{\mbox{22:~75,6}}
\end{pfaufg}
\pfgruppe{}
\begin{pfaufg}{P10 ’25}{\llap{\textasteriskcentered\,}23.}{}
Man beobachtet einen Bakterienstamm, der exponentiell wächst. Die Tabelle zeigt die Zahl der Bakterien in den ersten 4 Stunden. Zeige, dass der Wachstumsfaktor 1,4 ist. Gib eine Gleichung an, die die Zahl der Bakterien nach $t$ Stunden beschreibt. Prüfe mit einer Rechnung, ob diese Aussage stimmt: „Nach 24 Stunden sind es mehr als 250 000 Bakterien.“
\par\smallskip\noindent \sachtabelle{lrrrrr}{Zeit in Stunden & 0 & 1 & 2 & 3 & 4}{Anzahl der Bakterien & 80 & 112 & 157 & 220 & 308}\par
\fusshilfe{\mbox{23:~1,4; 80 · $1{,}4^{t}$; $\approx$ 257 136 > 250 000}}
\end{pfaufg}

\pfstufekopf[sgraphzuordnenundbegr]{Graph zuordnen und begründen}{in 1 der letzten 5 Prüfungen}
\pfgruppe{}
\begin{pfaufg}{}{24.}{}
Sieh dir im Koordinatensystem den Graphen $f$ an. Kreuze an, bei welchem Wert er die senkrechte Achse schneidet. 
\pfkreuzzeile{$0$}\pfkreuzzeile{$50$}\pfkreuzzeile{$100$}
\par\smallskip\noindent \begin{ksys}[xmin=0,xmax=10,ymin=0,ymax=700,xstep=1,ystep=100,xlabel=t in Jahren,ylabel=Wert,ablesen] \funktionab{100*1.2^\x}{f}{0}{10} \gerade{30}{100}{g} \parabel{6}{0}{0}{h} \end{ksys}\par
\fusshilfe{\mbox{24:~$100$}}
\end{pfaufg}
\pfgruppe{}
\begin{pfaufg}{P10 ’19}{\llap{\textasteriskcentered\,}25.}{}
Ein Ferkel wiegt 10 kg. Aus Erfahrung weiß man, dass seine Masse jede Woche um etwa 4 \% zunimmt. Zu Beginn jeder Woche wird es gewogen. Für jeden der drei Graphen: Entscheide, ob er zum Wachstum des Ferkels passt, und begründe.
\par\smallskip\noindent \begin{tabular}[t]{@{}c@{}}{\small 1}\enspace$\square$\\\begin{tikzpicture}\begin{axis}[width=3.9cm,height=3.4cm,scale only axis=false,axis lines=middle,xtick=\empty,ytick=\empty,xmin=0,xmax=1.08,ymin=0,ymax=1.08,clip=true,axis line style={-{Stealth[length=1.5mm]}},grid=both,minor tick num=0,xtick={0.2,0.4,0.6,0.8,1},ytick={0.2,0.4,0.6,0.8,1},xticklabels={},yticklabels={},grid style={mbgitter},xlabel={\tiny Zeit in Wochen},ylabel={\tiny Masse in kg},x label style={anchor=north east,at={(1,0)},yshift=-2pt},y label style={anchor=south west,at={(0,1)},xshift=1pt}]\addplot[black,line width=0.9pt,domain=0:1,samples=80] {0.25+0.6*x};\end{axis}\end{tikzpicture}\\[1pt]{\footnotesize $\square$ passt\quad $\square$ passt nicht}\\[2pt]\rule{3cm}{0.4pt}\end{tabular}\hspace{0.4cm}\begin{tabular}[t]{@{}c@{}}{\small 2}\enspace$\square$\\\begin{tikzpicture}\begin{axis}[width=3.9cm,height=3.4cm,scale only axis=false,axis lines=middle,xtick=\empty,ytick=\empty,xmin=0,xmax=1.08,ymin=0,ymax=1.08,clip=true,axis line style={-{Stealth[length=1.5mm]}},grid=both,minor tick num=0,xtick={0.2,0.4,0.6,0.8,1},ytick={0.2,0.4,0.6,0.8,1},xticklabels={},yticklabels={},grid style={mbgitter},xlabel={\tiny Zeit in Wochen},ylabel={\tiny Masse in kg},x label style={anchor=north east,at={(1,0)},yshift=-2pt},y label style={anchor=south west,at={(0,1)},xshift=1pt}]\addplot[black,line width=0.9pt,domain=0:1,samples=80] {0.2*exp(1.45*x)};\end{axis}\end{tikzpicture}\\[1pt]{\footnotesize $\square$ passt\quad $\square$ passt nicht}\\[2pt]\rule{3cm}{0.4pt}\end{tabular}\hspace{0.4cm}\begin{tabular}[t]{@{}c@{}}{\small 3}\enspace$\square$\\\begin{tikzpicture}\begin{axis}[width=3.9cm,height=3.4cm,scale only axis=false,axis lines=middle,xtick=\empty,ytick=\empty,xmin=0,xmax=1.08,ymin=0,ymax=1.08,clip=true,axis line style={-{Stealth[length=1.5mm]}},grid=both,minor tick num=0,xtick={0.2,0.4,0.6,0.8,1},ytick={0.2,0.4,0.6,0.8,1},xticklabels={},yticklabels={},grid style={mbgitter},xlabel={\tiny Zeit in Wochen},ylabel={\tiny Masse in kg},x label style={anchor=north east,at={(1,0)},yshift=-2pt},y label style={anchor=south west,at={(0,1)},xshift=1pt}]\addplot[black,line width=0.9pt,domain=0:1,samples=80] {0.88*(x^2)};\end{axis}\end{tikzpicture}\\[1pt]{\footnotesize $\square$ passt\quad $\square$ passt nicht}\\[2pt]\rule{3cm}{0.4pt}\end{tabular}\par
\fusshilfe{\mbox{25:~Graph 1 nein; Graph 2 ja; Graph 3 nein}}
\end{pfaufg}

\pfpruefstein{P10 ’26}
\begin{pfaufg}{}{}{}
Can bezieht Anfang 2026 eine neue Wohnung. Laut Mietvertrag steigt die Miete jedes Jahr um 1,9 \%. Im ersten Jahr zahlt er 650 € im Monat.
\end{pfaufg}
\begin{pfaufg}{}{a)}{2 BE}
Can legt eine Tabelle an. Vervollständige sie.
\par\smallskip\noindent \sachtabelle{lcccc}{Jahr & 2026 & 2027 & 2028 & 2029}{Miete in € & \leerzelle & 662,35 & 674,93 & \leerzelle}\par
\fusshilfe{\mbox{Prüfstein a):~2026: 650,00 €; 2029: 687,76 €}}
\end{pfaufg}
\begin{pfaufg}{}{b)}{3 BE}
Entscheide, welcher Graph zur Entwicklung der Miete passt, und kreuze ihn an. Erkläre bei den zwei übrigen Graphen jeweils, weshalb sie nicht passen.
\par\smallskip\noindent \begin{tabular}[t]{@{}c@{}}{\small A}\enspace$\square$\\\begin{tikzpicture}\begin{axis}[width=3.9cm,height=3.4cm,scale only axis=false,axis lines=middle,xtick=\empty,ytick=\empty,xmin=0,xmax=1.08,ymin=0,ymax=1.08,clip=true,axis line style={-{Stealth[length=1.5mm]}},grid=both,minor tick num=0,xtick={0.2,0.4,0.6,0.8,1},ytick={0.2,0.4,0.6,0.8,1},xticklabels={},yticklabels={},grid style={mbgitter},xlabel={\tiny Zeit in Jahren},ylabel={\tiny Miete in €},x label style={anchor=north east,at={(1,0)},yshift=-2pt},y label style={anchor=south west,at={(0,1)},xshift=1pt}]\addplot[black,line width=0.9pt,domain=0:1,samples=80] {0.88*(x^2)};\end{axis}\end{tikzpicture}\end{tabular}\hspace{0.4cm}\begin{tabular}[t]{@{}c@{}}{\small B}\enspace$\square$\\\begin{tikzpicture}\begin{axis}[width=3.9cm,height=3.4cm,scale only axis=false,axis lines=middle,xtick=\empty,ytick=\empty,xmin=0,xmax=1.08,ymin=0,ymax=1.08,clip=true,axis line style={-{Stealth[length=1.5mm]}},grid=both,minor tick num=0,xtick={0.2,0.4,0.6,0.8,1},ytick={0.2,0.4,0.6,0.8,1},xticklabels={},yticklabels={},grid style={mbgitter},xlabel={\tiny Zeit in Jahren},ylabel={\tiny Miete in €},x label style={anchor=north east,at={(1,0)},yshift=-2pt},y label style={anchor=south west,at={(0,1)},xshift=1pt}]\addplot[black,line width=0.9pt,domain=0:1,samples=80] {0.2*exp(1.45*x)};\end{axis}\end{tikzpicture}\end{tabular}\hspace{0.4cm}\begin{tabular}[t]{@{}c@{}}{\small C}\enspace$\square$\\\begin{tikzpicture}\begin{axis}[width=3.9cm,height=3.4cm,scale only axis=false,axis lines=middle,xtick=\empty,ytick=\empty,xmin=0,xmax=1.08,ymin=0,ymax=1.08,clip=true,axis line style={-{Stealth[length=1.5mm]}},grid=both,minor tick num=0,xtick={0.2,0.4,0.6,0.8,1},ytick={0.2,0.4,0.6,0.8,1},xticklabels={},yticklabels={},grid style={mbgitter},xlabel={\tiny Zeit in Jahren},ylabel={\tiny Miete in €},x label style={anchor=north east,at={(1,0)},yshift=-2pt},y label style={anchor=south west,at={(0,1)},xshift=1pt}]\addplot[black,line width=0.9pt,domain=0:1,samples=80] {0.25+0.6*x};\end{axis}\end{tikzpicture}\end{tabular}\par
\fusshilfe{\mbox{Prüfstein b):~B; A: Miete beginnt bei 0 € statt bei 650 €}}
\end{pfaufg}
\begin{pfaufg}{}{*c)}{3 BE}
Stelle eine Gleichung auf, mit der man die Monatsmiete nach beliebig vielen Jahren berechnen kann. Bestimme die Monatsmiete im Jahr 2040.
\rechenplatz{3}
\fusshilfe{\mbox{Prüfstein c):~650 · $1{,}019^{t}$; $\approx$ 845,96 €}}
\end{pfaufg}
\end{document}